% 第二部：篇名與全文架構。出處：課程倉 1151/seminar/my/導讀_W06_Waldron_細目次與大綱.md 第二至四節；翻譯倉附錄圖「全文的論證結構」。
\chapter{篇名與全文架構}

\qaentry{第二部之一}{篇名的 the core of the case 與正文的 core case}

\begin{qaquestion}
篇名的 the core of the case 指什麼？正文另有 core case 一詞，兩者是否同義？
\end{qaquestion}

篇名的 the core of the case 指反對違憲審查的論據中，抽離各國歷史與個案結果之後所餘的部分\pg{1351}，譯本譯作「反對違憲審查的核心論據」。正文的 core case 指四項假設全部成立的社會\pg{1359--1360}，譯作「核心情形」，與第柒部分的 non-core cases（非核心情形）成對。第參部分末段把兩者接在一起：結果面的理由至多未有定論，程序面的理由幾乎全部不利於違憲審查，立法機關一方在兩種問法下都勝出，作者說 \en{that will be the core of the case against judicial review}\pg{1375--1376}。

\qaentry{第二部之二}{七個部分的排序}

\begin{qaquestion}
本文的七個部分各回答什麼問題？排序本身構成怎樣的論證？
\end{qaquestion}

\begin{qafigure}{全文的論證結構（1353--1406）。改繪自譯本附錄〈全文的論證結構〉。\label{fig:structure}}
% 全文的論證結構，改繪自翻譯倉 JudicialReview/Werke/Teile/93_Argument_Figures.tex。灰階為底，結論用強調色。
\begin{tikzpicture}[qa,node distance=4.5mm]
  \newcommand{\p}[1]{\\[-1pt]{\scriptsize\color{dimink}#1}}
  \node[gbox,text width=30mm] (a1) at (0,0) {一、民主制度\p{1361--1362}};
  \node[gbox,text width=30mm,right=2mm of a1] (a2) {二、司法制度\p{1363--1364}};
  \node[gbox,text width=30mm,right=2mm of a2] (a3) {三、對權利的承諾\p{1364--1366}};
  \node[gbox,text width=30mm,right=2mm of a3] (a4) {四、關於權利的分歧\p{1366--1369}};
  \node[hdr,anchor=south west] at ([yshift=1.2mm]a1.north west) {貳、四項假設};
  \coordinate (bar) at ([yshift=-3.5mm]a1.south);
  \node[nbox,text width=64mm,anchor=north] (core) at ([yshift=-4mm,xshift=24mm]bar) {核心情形：四項假設均成立\p{1359--1360}};
  \node[nbox,text width=80mm,below=of core] (form) {參、論證的形式：權利分歧須經決定程序解決；評價程序的理由分為結果面與程序面\p{1369--1376}};
  \node[nbox,text width=38mm,below=of form.south west,anchor=north west] (out) {肆、結果面的理由：個案取向、文本取向、說明理由；至多未有定論\p{1376--1386}};
  \node[nbox,text width=38mm,below=of form.south east,anchor=north east] (proc) {伍、程序面的理由：由誰決定、以何種程序決定；大多有利於立法機關\p{1386--1395}};
  \node[nbox,text width=80mm,anchor=north] (tyr) at ([yshift=-4.5mm]out.south -| form) {陸、多數暴政：在第三、四項假設之下，不應預期議題上的少數與決定上的少數重合\p{1395--1401}};
  \node[kbox,text width=80mm,below=of tyr] (con) {結論：權利分歧應以立法制度解決，以權利為本的違憲審查不適合這樣的社會\p{1360；1406}};
  \node[nbox,text width=31mm,anchor=north west] (nc) at ([xshift=5mm]form.north east) {柒、非核心情形：民主制度失能、成員甚少關心他人權利、孤立而隔絕的少數；反對論不成立，違憲審查也不因此成立\p{1401--1406}};
  \foreach \x in {a1,a2,a3,a4} \draw[figGray,line width=0.5pt] (\x.south) -- (\x.south |- bar);
  \draw[figGray,line width=0.5pt] (a1.south |- bar) -- (a4.south |- bar);
  \draw[arr] (bar -| core) -- (core);
  \draw[arr] (core) -- (form);
  \draw[arr] (form.south -| out) -- (out);
  \draw[arr] (form.south -| proc) -- (proc);
  \draw[arr] (out.south) -- (out.south |- tyr.north);
  \draw[arr] (proc.south) -- (proc.south |- tyr.north);
  \draw[arr] (tyr) -- (con);
  \draw[dsh] (bar -| nc) -- node[lbl,fill=white,right] {假設不成立} (nc.north);
  \draw[dsh] (tyr.east) -| node[lbl,fill=white,pos=0.75,right] {兩種少數普遍重合} (nc.south);
\end{tikzpicture}

\end{qafigure}

導論末段已列出同一順序\pg{1353}。依各部分回答的問題，排序如下。

\begin{longtable}{@{}L{0.25\linewidth}L{0.30\linewidth}L{0.37\linewidth}@{}}
\toprule
部分 & 要回答的問題 & 作用 \\
\midrule
\endhead
壹、違憲審查的界定\newline\textcolor{dimink}{\small 1353--1359} & 攻擊的是哪一種違憲審查 & 標的：限於對立法的強式審查 \\
貳、四項假設\newline\textcolor{dimink}{\small 1359--1369} & 在什麼樣的社會裡討論 & 前提：四項假設界定核心情形 \\
參、論證的形式\newline\textcolor{dimink}{\small 1369--1376} & 以什麼標準選擇決定程序 & 判準：結果面與程序面兩類理由 \\
肆、結果面的理由\newline\textcolor{dimink}{\small 1376--1386} & 法院是否較可能得出正確的權利決定 & 檢驗一：法院被歸予的三項優點 \\
伍、程序面的理由\newline\textcolor{dimink}{\small 1386--1395} & 撇開結果，誰有權作終局決定 & 檢驗二：向不同意者說明的兩個問題 \\
陸、多數暴政\newline\textcolor{dimink}{\small 1395--1401} & 多數暴政之說能否支持違憲審查 & 對最常見反駁的回應 \\
柒、非核心情形\newline\textcolor{dimink}{\small 1401--1406} & 假設不成立時如何 & 前提不成立時論證的效力 \\
\bottomrule
\end{longtable}

第貳部分的四項假設各有後文取用之處：B 節的法院特徵在第肆部分 A 節與第伍部分取用，C 節在第陸、柒部分取用，D 節在第參部分與第肆部分 B 節取用，A 節在第伍、柒部分取用。非核心情形在第柒部分之前已多次預告\pg{1360、1366、1377、1386、1398}。

\qaentry{第二部之三}{小標未顯示的內容}

\begin{qaquestion}
刊本只有兩層共 16 條小標。導讀時有哪些內容在小標上看不出來，應向聽者指出？
\end{qaquestion}

\begin{enumerate}
\item 第壹部分只有一個小標，內容卻有四項區分：強式與弱式、以權利為對象的審查與一般憲法審查、事後的個案審查與事前的抽象審查、普通法院與專設的憲法法院；後三項各帶作者的評論\pg{1354--1359}。
\item 第伍部分後半回應五項調和違憲審查與民主的論證，小標未列\pg{1393--1395}。
\item 第陸部分末尾回應 Dworkin 的「背景條件」說與 nemo iudex in sua causa（任何人不得為自己案件之法官）\pg{1398--1401}。
\end{enumerate}
